\centerline{\bf \Huge Гомотетия}
\centerline{\bf \Huge Level: Hard}

\begin{flushright}
        \scriptsize{это уже не гомотетия\\ это homothety}
\end{flushright}

\problem{1}
Окружности $\omega$ и $\omega_a$ с центрами в точках $I$ и $I_a$
являются вписанной и вневписанной напротив вершины $A$ окружностями
треугольника $ABC$ соответственно. Эти окружности касаются стороны $BC$
в точках $A_1$ и $A_2$ соответственно. Также $AH$ — высота треугольника
$ABC$, опущенная из вершины $A$.\

(a) Докажите, что одна из точек пересечения прямой $AA_2$ с окружностью
$\omega$ диаметрально противоположна точке $A_1$.\

(b) Докажите, что $\angle IHC=\angle I_aHC$.\

(c) Докажите, что точка Нагеля $N$, точка $I$ и точка пересечения медиан $M$
лежат на одной прямой.

\problem{2}
В треугольник $ABC$ вписана окружность $\omega$. Точки $M_a$, $M_b$ и
$M_c$ — середины сторон $BC$, $CA$ и $AB$ соответственно. Точки $K_a$,
$K_b$ и $K_c$ симметричны точкам касания окружности $\omega$ со
сторонами $BC$, $CA$ и $AB$ относительно биссектрис углов $\angle A$,
$\angle B$ и $\angle C$ соотвественно. Докажите, что прямые $M_aK_a$,
$M_bK_b$ и $M_cK_c$ пересекаются в одной точке.

\problem{3}
Окружности $\omega_1$, $\omega_2$ и $\omega_3$ попарно касаются друг
друга внешним образом в точках $A$, $B$ и $C$. На окружности $\omega_1$
взята произвольная точка $P$. Прямая $AP$ вторично пересекает окружность
$\omega_2$ в точке $Q$, прямая $QB$ вторично пересекает окружность
$\omega_3$ в точке $R$, а прямая $RC$ вторично пересекает окружность
$\omega_1$ в точке $S$. Докажите, что точки $P$ и $S$ диаметрально
противоположны.

\problem{4}
В треугольнике $ABC$ проведена биссектриса $AD$. Общая касательная
$\ell$ к окружностям $(ABD)$ и $(ACD)$ касается их в точках $M$ и $N$
соответственно. Докажите, что прямая $\ell$ касается окружности,
проходящей через середины отрезков $MN$, $BD$, $CD$.

\problem{5}
В остроугольном треугольнике $ABC$ проведена высота $AA_1$ и отмечена
точка $M$ пересечения медиан. Луч $MA_1$ пересекает окружность $(ABC)$ в
точке $K$. Докажите, что окружность $(A_1BK)$ касается прямой $AB$.


\newpage

\hspace{3mm}

\problem{6}
В треугольник $ABC$ вписана окружность $\omega$ с центром в точке $I$,
касающаяся стороны $BC$ в точке $D$. Пусть $\omega_B$ и $\omega_C$ —
вписанные окружности треугольников $ABD$ и $ACD$ соответственно,
касающиеся стороны $BC$ в точках $E$ и $F$ соответственно. Пусть $P$ —
точка пересечения отрезка $AD$ с линией центров окружностей $\omega_B$ и
$\omega_C$, $X$ — точка пересечения прямых $BI$ и $CP$, $Y$ — точка
пересечения прямых $CI$ и $BP$. Докажите, что прямые $EX$ и $FY$
пересекаются на окружности $\omega$.

\problem{7}
Две окружности $\Omega$ и $\omega$ касаются в точке $A$. Прямая $\ell$
пересекает окружность $\Omega$ в точках $B$ и $C$ и касается окружности
$\omega$ в точке $D$. Докажите, что прямая $AD$ проходит через середину
дуги окружности $\Omega$, стягиваемую хордой $BC$ (разберите случаи
внешнего и внутреннего касания).\

\problem{8}
Вокруг равнобедренного треугольника $ABC$ ($AB=AC$) описана окружность
$\Omega$. Окружность $\omega$ касается сторон $AB$ и $AC$ в точках $K$ и
$L$ и касается внутренним образом окружности $\Omega$. 

(a) Докажите, что центр вписанной окружности треугольника $ABC$ лежит на отрезке $KL$.\

(b) Докажите утверждение предыдущего пункта для произвольного треугольника
$ABC$.